\section{A sufficient two-opportunity mechanism}\label{app:canonical-mechanism}

The construction in Section~\ref{sec:model-observation} uses $N\geq4$ firms,
$b,v,Q,\nu>0$, distinct shares $0<p_i<1$ with $\sum_i p_i=N/2$, and $D=b+NQ/2>2Q$.
Baseline capacities are $(Qp_i,Q(1-p_i))$.
At each node, maintain the additional restriction that the supporting linear
cost $c_i=v(D-q_i)/D^2-\nu q_i$ is nonnegative, and choose quadratic cost
$\nu q_i^2/2$, so the baseline marginal payoff vanishes.
The following derivation concerns this interior origin and its tangent
response, not unrestricted finite shocks.

For one node with baseline capacities $q_i$, differentiating the
marginal-cost-less-marginal-contribution field gives diagonal entry
$\nu+2v(D-q_i)/D^3$ and off-diagonal row entry $v(D-2q_i)/D^3$.
Consequently
\begin{equation}\label{eq:canonical-jacobian}
  J=dI+u\mathbf1^\top,\quad d=\nu+v/D^2,\quad
  u_i=v(D-2q_i)/D^3>0,\quad f_i=(D-q_i)/D^2.
\end{equation}
Here $f_i$ is the opportunity forcing.
Writing $U=\sum_i u_i=Nv(D-Q)/D^3$ and $F=\sum_i f_i=N(D-Q/2)/D^2$, define $d_A=d+U$.
The projectors $P_A=u\mathbf1^\top/U$ and $P_\perp=I-P_A$ obey $JP_A=d_A P_A$
and $JP_\perp=dP_\perp$.
Applying them to the anchored recurrence yields persistence modes
$a=\kappa/(d+\kappa)$ and $a_A=\kappa/(d_A+\kappa)$.
For a permanent unit node improvement, the capacity response after $h$ stages is
\begin{equation}\label{eq:canonical-step}
  x_h(q)=\frac{f(q)-(F/U)u(q)}{d}(1-a^h)
  +\frac{(F/U)u(q)}{d_A}(1-a_A^h),\qquad x_0(q)=0.
\end{equation}
This follows by summing the two geometric recurrences, with $0\leq a_A\leq a<1$.
The aggregate node response is $T_h=F(1-a_A^h)/d_A$.

The response $x_h(q)$ is affine in $q$ and positive on $[0,Q]$ for $h\geq1$.
At $q=Q$, both modal coefficients in \eqref{eq:canonical-step} are positive, since
$f(q)-(F/U)u(q)=(q-Q/2)/[D(D-Q)]$.
At $q=0$, write the response as
$A(1-a_A^h)-B(1-a^h)$, with
$A=(D-Q/2)/[D(D-Q)d_A]$ and $B=Q/[2D(D-Q)d]$.
Both constants are positive, and $A>B$ because
$2(D-Q)d-QU=2(D-Q)(\nu+vb/D^3)>0$.
Since $1-a_A^h\geq1-a^h>0$, this endpoint response is also positive; affinity
covers the interval.
A relative innovation raises the plus-node capacity by $x_h(Qp_i)$ and lowers
the minus-node capacity by $x_h(Q(1-p_i))$.
Differentiating the normalized plus share therefore gives
\begin{equation}\label{eq:canonical-share-step}
  \xi_i(h)=\frac{(1-p_i)x_h(Qp_i)+p_i x_h(Q(1-p_i))}{Q}>0
  \quad(h\geq1).
\end{equation}
The positive share response is generated by the joint game, including both modes.

The static envelope evaluates cash flow along exactly this capacity response.
Affinity of $x_h(q)$ and the balanced baseline shares imply that firm $i$'s
payout derivative is $(2p_i-1)L_h$, where
\begin{equation}\label{eq:canonical-payout-scale}
  L_h=\frac QD+\frac{vQ}{D^2}[x_h(Q)-T_h].
\end{equation}
This scale is strictly positive.
Indeed, $x_h(Q)\geq0$, $T_h\leq F/d_A$, and
$Dd_A-vF=D\nu+vb/D^2>0$, which bounds $L_h$ below by a positive constant.
The bounded stable responses make $L_\delta=\sum_{h\geq0}\delta^h L_h$ finite
and positive for $0<\delta<1$.
Substituting the supporting costs into baseline cash flow gives
$\bar\pi_i=B[1+(2p_i-1)^2]$, with $B=(v/D^2+\nu/2)Q^2/2>0$.
Dividing discounted payout news by the positive ex-payout price yields
\begin{equation}\label{eq:canonical-return-shape}
  g(p_i)=G\frac{2p_i-1}{1+(2p_i-1)^2},\qquad
  G=\frac{(1-\delta)L_\delta}{\delta B}>0.
\end{equation}
Its derivative is $2G[1-(2p-1)^2]/[1+(2p-1)^2]^2>0$ on $0<p<1$.
Thus the financial ordering follows from the same joint operating response
and constant-discount valuation convention, rather than an independently
imposed return ranking.

For the observation laws in \eqref{eq:axis-observation}, common mass can stay
at each antipodal atom; surplus component mass must move to the other axis at
squared cost two per unit.
This constructs cost $2|p-q|$, and the required surplus provides the matching
lower bound for every coupling.
At distinct baseline shares, approach toward reference $j$ has derivative
$2o_{ij}$ times the target share change, with $o_{ij}=\operatorname{sgn}(p_j-p_i)$.
Pooling dated probability laws before evaluating transport preserves the
corresponding local weighted share derivative.

Let $f_{iys}\geq0$ be the article-timing weight of source date $s$ and
$k_{iys}\geq0$ the later-year article-weighted response assembled from
\eqref{eq:canonical-share-step}.
For sources in year $y$, the earlier article window contains no observations
after the source; strict timing excludes same-date financial news.
With centered independent primitive innovations $\varepsilon_s$, finite
variances $\sigma_s^2$, and no cross moments with initial gaps or
proxy/observation errors, the tangent centered moment is
\begin{equation}\label{eq:canonical-directional-moment}
  E[\operatorname{Cov}_{j\ne i}(R_{ijy},A_{ijy})]
  =2\operatorname{Cov}_{j\ne i}(g(p_j),o_{ij})
  \sum_{s\in y}f_{iys}k_{iys}\sigma_s^2\geq0.
\end{equation}
To see the last sign, express the reference covariance as half the equally
weighted ordered pair sum of
$[g(p_j)-g(p_k)](o_{ij}-o_{ik})$.
Both factors have the same ordering in baseline share, so every product is nonnegative.
Independence across source dates removes cross-source expectations; primitive
news can be shared by every firm.
Nonnegative pooling across target-years preserves the sign.
General response and financial loading matrices need not have this ordering,
and observed raw returns and reports need not satisfy the exclusion restrictions.
The aligned operating and return responses therefore give a sufficient
local mechanism for the expected moment in
Proposition~\ref{prop:directional-alignment}.
